\documentclass[10pt,a4paper]{amsart}
\usepackage[margin=19mm]{geometry}
\usepackage{amsmath,amssymb,mathtools}
\usepackage[scr=rsfs,cal=euler]{mathalfa}
\usepackage{xfrac}
\usepackage[unicode,hidelinks,bookmarks=false]{hyperref}
\newcommand{\ie}{i.e.\ }
\newcommand{\cf}{cf.\ }
\newcommand{\resp}{respectively\ }
\newcommand{\Subsection}{\subsection}
\providecommand{\nobreakdash}{\nobreak\hbox{-}}
\setlength{\emergencystretch}{2em}
\multlinegap0pt
\usepackage{amssymb}
\usepackage{mathtools}
\newtheorem{lemma}{Lemma}
\newcommand{\R}{\mathbb R}
\newcommand{\Z}{\mathbb Z}
\title{A simple derivation of the zero-temperature BCS functional from the reduced BCS Hamiltonian}
\author{Christian Hainzl, Riccardo Panza}
\date{Source: arXiv:2609.04060v1 (CC BY 4.0)}
\begin{document}
\maketitle

\noindent\emph{Source and licence.} A simple derivation of the zero-temperature BCS functional from the reduced BCS Hamiltonian. Version arXiv:2609.04060v1 (CC BY 4.0), distributed by arXiv under the Creative Commons Attribution 4.0 International licence, http://creativecommons.org/licenses/by/4.0/, as recorded in the arXiv licence metadata for this exact version and shown on the abstract page https://arxiv.org/abs/2609.04060v1. Full licence notice: \texttt{licenses/arxiv-2609.04060-CC-BY-4.0.txt}.\par\smallskip

\noindent\textit{Editorial scope.} Counted unit: the article's lemma lem:abstract-localization (source0.tex lines 623--639). The article's complete original proof of it is delivered with it (source0.tex lines 641--791, one \texttt{proof} environment of 151 lines). No mathematics is written, repaired or re-derived: the statement and the proof are the article's own lines, in the article's own notation, and no retained line is substituted or re-flowed.  Retained uncounted as the source-local context the proof needs: lines 613--617.  Nothing was omitted from any retained span, so the package carries no omission record.  No retained line contains a citation, so the wrapper carries no bibliography and no bibliography entry.  No retained line contains a figure environment, an image-inclusion command or drawing code, so no image is delivered and the package declares no asset dependency.  Editorial formatting only, in addition: an AMS \texttt{amsart} preamble that restates the article's own declarations and macros used by the retained lines, namely the environments \texttt{lemma} on separate counters, together with the standard packages those lines need.  The retained environments print in this document as Lemma 1 in order of appearance, whereas the article prints the same items with the numbering of its own full document.  The complete native source of the article as distributed in the arXiv e-print for this exact version is bundled unchanged as \texttt{sources/209341/source0.tex}.
\begin{equation}\label{eq:chi-theta}
 \chi_j=\chi\left(\frac{X-j\ell}{\ell}\right),\qquad
 \theta_k=\chi\left(\frac{Y-k\ell}{\ell}\right),\qquad
 A_{jk}=\chi_j\theta_k.
\end{equation}

\begin{lemma}[Two-variable localization]\label{lem:abstract-localization}
Using the above definitions \eqref{eq:chi-theta}, there is a constant $C_\chi<\infty$ (depending only on the function $\chi$) such that
\begin{equation}\label{eq:IMS-general}
 H_{L, \mu}=\sum_{j,k}A_{jk}^*H_{L, \mu}A_{jk}+\mathcal{E}_{\rm loc}
\end{equation}
as quadratic forms, where
\begin{equation}\label{eq:Eloc-bound-general}
 \|\mathcal{E}_{\rm loc}\|
 \le \frac{C_\chi}{\ell^2}
 \bigl(\|[X,[H_{L, \mu},X]]\|+\|[Y,[H_{L, \mu}, Y]]\|\bigr).
\end{equation}
Moreover, for every $\psi \in \mathcal{F}_f$:
\begin{equation}\label{eq:Zlocal-general}
 \sum_{j,k \in \Z}\|(X+iY-(j+ik)\ell)A_{jk}\psi\|^2
 \le C_\chi\left(\ell^2+\frac{\|[X,Y]\|^2}{\ell^2}\right)\|\psi\|^2.
\end{equation}
\end{lemma}

\begin{proof}
We apply an IMS-type identity successively to the two spectral localizations. For a bounded self-adjoint partition of unity satisfying $\sum_j\chi_j^2=1$ and a self-adjoint operator $H$ bounded from below, one has
\[
 \sum_j\chi_j H\chi_j - H
 =
 \frac12\sum_j[\chi_j,[H,\chi_j]]
\]
as a quadratic-form identity. Applying this twice, we find
\begin{align*}
   \sum_{j,k}A_{jk}^*HA_{jk}
 &=\sum_k\theta_k\left(\sum_j\chi_jH\chi_j\right)\theta_k \\
 &=\sum_k\theta_k \left(H + \frac{1}{2}\sum_{j}\left[\chi_j,[H, \chi_j]\right]\right)\theta_k = \\
 &= H + \frac{1}{2}\sum_{k}\left[\theta_k,[H, \theta_k]\right] + \frac{1}{2}\sum_{j,k} \theta_k\left[\chi_j,[H, \chi_j]\right] \theta_k =: H - \mathcal{E}_{loc}\;.
\end{align*}
We first estimate the contribution arising from the localization in the $Y$-variable, namely $\sum_k [\theta_k,[H,\theta_k]]$.
Using the Fourier representation of $\chi$, we write
\[
 \chi(t)=\frac1{\sqrt{2\pi}}
 \int_{\R}\widehat\chi(\xi)e^{i \xi t}\,d\xi \quad\Longrightarrow\quad \theta_k
 =
 \frac1{\sqrt{2\pi}}
 \int_{\R}\widehat\chi(\xi) e^{-i \xi k}e^{i \xi Y/\ell}\,d\xi\;,
\]
so that
\[
\sum_{k \in \Z} [\theta_k,[H,\theta_k]] = \frac{1}{2\pi}\sum_{k \in \mathbb{Z}}\int_{\mathbb{R}\times\mathbb{R}}d\xi\,dw\,e^{-i(\xi + w) k}\,\hat{\chi}(\xi)\hat{\chi}(w)  \left[e^{i\xi Y/\ell}, \left[H, e^{i w Y/\ell}\right]  \right] .
\]

Summing over the localization index is then handled by the Poisson summation formula
\[
 \sum_{k\in\Z}e^{-i(\xi+w)k}
 =
 2\pi\sum_{m\in\Z}\delta(\xi+w-2\pi m)
\]
to get
\[
\sum_{k \in \Z} [\theta_k,[H,\theta_k]] = \sum_{m \in \Z}\int_{\R}dw\,\hat{\chi}(w)\,\hat{\chi}(2\pi m - w)\left[e^{i(2 \pi m - w) Y / \ell}, \left[H, e^{i w Y/ \ell}\right]  \right].
\]
The remaining double commutator is estimated by a Duhamel expansion. Applying the fundamental theorem of calculus to the corresponding unitary conjugation $f(r) := e^{i(1-r)tY/\ell}He^{irtY/\ell}$ yields
\[
 [H,e^{itY/\ell}]
 =
 \frac{it}{\ell}
 \int_0^1
 e^{i(1-r)tY/\ell}[H,Y]e^{irtY/\ell}\,dr\,,
\]
and a second application gives the bound
\[
 \|[e^{isY/\ell},[H,e^{itY/\ell}]]\|
 \le
 \frac{|st|}{\ell^2}\|[Y,[H,Y]]\|\,.
\]

Therefore
\begin{align*}
 \left\|\sum_k[\theta_k,[H,\theta_k]]\right\|
 &\le
 \frac{\|[Y,[H,Y]]\|}{\ell^2}
 \sum_{m\in\Z}\int_{\R}
 |s|\,|2\pi m-s|\,
 |\widehat\chi(s)|\,|\widehat\chi(2\pi m-s)|\,ds\\
 &=\frac{C_{\chi}}{\ell^2}
 \|[Y,[H,Y]]\|.
\end{align*}
The constant $C_\chi$ is finite because $\widehat\chi$ is a Schwartz function. The contribution from the localization in the $X$-variable is estimated in exactly the same way. This proves the bound on $\mathcal E_{\rm loc}$.

It remains to prove the localization estimate \eqref{eq:Zlocal-general}. We treat the $X$- and $Y$-components separately.
Since, by functional calculus, $(X-j\ell)$ commutes with $\chi_j$, we get 
\begin{align*}
 &\sum_{j,k}\|(X-j\ell)\chi_j\theta_k\psi\|^2\\
 &\quad=
 \sum_k
 \left\langle
 \theta_k\psi,
 \left[
 \sum_j\chi_j\,(X-j\ell)^2\,\chi_j
 \right]
 \theta_k\psi
 \right\rangle.
\end{align*}
Introduce the periodic function
\[
 h_\chi(t):=
 \sum_{j\in\mathbb Z}(t-j)^2\chi(t-j)^2 .
\]
The sum is locally finite, so \(h_\chi\) is continuous and in particular
\(
 C_\chi:=\sup_{r\in\mathbb R}h_\chi(r)<\infty.
\)
By functional calculus,
\[
 \sum_j\chi_j(X-j\ell)^2\chi_j
 =
 \ell^2h_\chi(X/\ell)
 \le C_\chi\ell^2.
\]
Consequently,
\begin{equation*}
 \sum_{j,k}\|(X-j\ell)A_{jk}\psi\|^2
 \le 
 C_\chi\ell^2\|\psi\|^2.
 \label{eq:X-local-piece}
\end{equation*}

The $Y$-component requires one additional commutator, because $\chi_j$ is a function of $X$ rather than of $Y$. We write
\[
 (Y-k\ell)\chi_j\theta_k
 =
 \chi_j(Y-k\ell)\theta_k+[Y,\chi_j]\theta_k.
\]
The first term is estimated exactly as in the $X$-component. It therefore remains only to control the commutator term. We claim that
\begin{equation}\label{eq:vector-comm-est}
 \left\|\sum_j[Y,\chi_j]^*[Y,\chi_j]\right\|
 \le \frac{C_\chi}{\ell^2}\|[X,Y]\|^2.
\end{equation}
To prove this estimate, we again use the Fourier representation of $\chi$ and sum over $j$ with the Poisson summation formula. This gives
\[
 \sum_j[Y,\chi_j]^*[Y,\chi_j] = \sum_{m \in \Z}\int_{\R}ds\,\hat{\chi}(s)\overline{\hat{\chi}(s + 2\pi m)}\left[Y, e^{i(2\pi m + s) X/\ell}\right]^* \left[Y, e^{i s X/ \ell}\right]\;.
\]
The same Duhamel argument as above yields
\[
 \|\left[Y,e^{isX/\ell}\right]\|
 \le \frac{|s|}{\ell}\|[X,Y]\|.
\]
Hence
\begin{align*}
 \left\|\sum_j[Y,\chi_j]^*[Y,\chi_j]\right\|
 &\le \frac{\|[X,Y]\|^2}{\ell^2}
 \sum_{m\in\Z}\int_{\R}
 |s|\,| s +2\pi m|\,
 |\widehat\chi(s)|\,|\widehat\chi(s + 2\pi m )|\,ds\\
 &=\frac{C_\chi}{\ell^2}\|[X,Y]\|^2.
\end{align*}
This proves \eqref{eq:vector-comm-est}. Consequently,
\begin{align*}
 \sum_{j,k}\|[Y,\chi_j]\theta_k\psi\|^2
 &=
 \sum_k
 \left\langle
 \theta_k\psi,
 \left(\sum_j[Y,\chi_j]^*[Y,\chi_j]\right)
 \theta_k\psi
 \right\rangle\\
 &\le
 \frac{C_\chi}{\ell^2}\|[X,Y]\|^2
 \sum_k\|\theta_k\psi\|^2\\
 &=
 \frac{C_\chi}{\ell^2}\|[X,Y]\|^2\|\psi\|^2.
\end{align*}
Combining this commutator estimate with the preceding bounds for the $X$- and $Y$-components proves \eqref{eq:Zlocal-general} and completes the proof.
\end{proof}

\end{document}
